\documentclass[12pt,a4paper]{report}
\usepackage[utf8]{inputenc}
\usepackage{graphicx}
\usepackage{listings}
\usepackage{xcolor}
\usepackage{geometry}
\usepackage{hyperref}
\geometry{a4paper, margin=1in}

\definecolor{codegreen}{rgb}{0,0.6,0}
\definecolor{codegray}{rgb}{0.5,0.5,0.5}
\definecolor{codepurple}{rgb}{0.58,0,0.82}
\definecolor{backcolour}{rgb}{0.95,0.95,0.92}

\lstdefinestyle{mystyle}{
    backgroundcolor=\color{backcolour},   
    commentstyle=\color{codegreen},
    keywordstyle=\color{magenta},
    numberstyle=\tiny\color{codegray},
    stringstyle=\color{codepurple},
    basicstyle=\ttfamily\footnotesize,
    breakatwhitespace=false,         
    breaklines=true,                 
    captionpos=b,                    
    keepspaces=true,                 
    numbers=left,                    
    numbersep=5pt,                  
    showspaces=false,                
    showstringspaces=false,
    showtabs=false,                  
    tabsize=2
}
\lstset{style=mystyle}

\title{
    \vspace{2cm}
    \Huge \textbf{MULTI-TIER ATTENDANCE HUB} \\
    \vspace{1cm}
    \Large A Comprehensive Academic Project Report \\
    \vspace{0.5cm}
    \large (Module 1 \& Module 2)
}
\author{System Architecture \& Implementation Report}
\date{\today}

\begin{document}

\maketitle
\tableofcontents
\newpage

\chapter{Problem Identification, Requirement Engineering \& System Design Phase}

\section{Problem Identification \& Feasibility Study}

\subsection{Identification of a Real-World Problem}
Traditional methods of taking attendance in educational institutions—such as calling out roll numbers or passing a sign-in sheet—are highly inefficient, time-consuming, and susceptible to proxy attendance. For large classes, manual attendance consumes a significant portion of instructional time. The need for an automated biometric system is paramount to maintaining academic integrity and maximizing teaching hours.
Biometrics are body measurements and calculations related to human characteristics and features. Biometric authentication (or realistic authentication) is used in computer science as a form of identification and access control. It is also used to identify individuals in groups that are under surveillance.
Biometric identifiers are the distinctive, measurable characteristics used to label and describe individuals. Biometric identifiers are often categorized as physiological characteristics which are related to the shape of the body. Examples include, but are not limited to fingerprint, palm veins, face recognition, DNA, palm print, hand geometry, iris recognition, retina, odor/scent, voice and shape of ears. Behavioral characteristics are related to the pattern of behavior of a person, including but not limited to mouse movement, typing rhythm, gait, signature, voice, and behavioral profiling. Some researchers have coined the term "Behaviometrics" (behavioral biometrics) to describe the latter class of biometrics.
More traditional means of access control include token-based identification systems, such as a driver's license or passport, and knowledge-based identification systems, such as a password or personal identification number. Since biometric identifiers are unique to individuals, they are more reliable in verifying identity than token and knowledge-based methods; however, the collection of biometric identifiers raises privacy concerns. 


== Biometric functionality ==
Many different aspects of human physiology, chemistry or behavior can be used for biometric authentication. The selection of a particular biometric for use in a specific application involves a weighting of several factors. Jain et al. (1999) identified seven such factors to be used when assessing the suitability of any trait for use in biometric authentication. Biometric authentication is based upon biometric recognition which is an advanced method of recognizing biological and behavioral characteristics of an Individual.
Factors include:

Universality: Every person using a system should possess the trait.
Uniqueness: The trait should be sufficiently different for individuals in the relevant population such that they can be distinguished from one another.
Permanence: The manner in which a trait varies over time. More specifically, a trait with good permanence will be reasonably invariant over time with respect to the specific matching algorithm.
Measurability (collectability): The ease of acquisition or measurement of the trait. In addition, acquired data should be in a form that permits subsequent processing and extraction of the relevant feature sets.
Performance: The accuracy, speed, and robustness of technology used (see performance section for more details).
Acceptability: How well individuals in the relevant population accept the technology such that they are willing to have their biometric trait captured and assessed.
Circumvention: The ease with which a trait might be imitated using an artifact or substitute.
Proper biometric use is application dependent. Certain biometrics will be more suitable than others based on the required levels of convenience and security. No single biometric will meet all the requirements of every possible application.

The block diagram illustrates the two basic modes of a biometric system. First, in verification (or authentication) mode the system performs a one-to-one comparison of a captured biometric with a specific template sto



\subsection{Problem Justification}
An automated biometric attendance system is justified as it reclaims instructional time, eliminates human error, and completely eradicates attendance fraud. Utilizing facial recognition provides a frictionless experience without requiring specialized hardware like fingerprint scanners, as it utilizes the existing ubiquitous web cameras on teacher and student devices.
Facial recognition systems are applications that use algorithms to compare facial images or representations to verify an identity or to identify a person from a set of known identities. One-to-one verification compares a facial image to a reference image that is linked to a claimed identity. One-to-many identification is a search of an image against a database to determine if the image matches a database ID.
Development on similar systems began in the 1960s as a form of computer application. Since their inception, facial recognition systems have seen wider uses in recent times on smartphones and in other forms of technology, such as robotics. Because computerized facial recognition involves the measurement of a human's physiological characteristics, facial recognition systems are categorized as biometrics. Although the accuracy of facial recognition systems as a biometric technology is lower than iris recognition, fingerprint image acquisition, palm recognition or voice recognition, it is widely adopted due to its contactless process. Facial recognition systems have been deployed in advanced human–computer interaction, video surveillance, law enforcement, passenger screening, decisions on employment and housing, and automatic indexing of images.
Facial recognition systems are employed throughout the world today by governments and private companies. Their effectiveness varies, and some systems have previously been scrapped because of their ineffectiveness. The use of facial recognition systems has also raised controversy, with claims that the systems violate citizens' privacy, commonly make incorrect identifications, encourage gender norms and racial profiling, and do not protect important biometric data. The appearance of synthetic media such as deepfakes has also raised concerns about its security. These claims have led to the ban of facial recognition systems in several cities in the United States. Growing societal concerns led social networking company Meta Platforms to shut down its Facebook facial recognition system in 2021, deleting the face-scan data of more than one billion users. The change represented one of the largest shifts in facial recognition usage in the technology's history. IBM also stopped offering facial recognition technology due to similar concerns.


== History of facial recognition technology ==


=== Initial development ===
Automated facial recognition was pioneered in the 1960s by Woody Bledsoe, Helen Chan Wolf, and Charles Bisson, whose work focused on teaching computers to recognize human faces. Their early facial recognition project was dubbed "man-machine" because a human first needed to establish the coordinates of facial features in a photograph before they could be used by a computer for recognition. Using a graphics tablet, a human would pinpoint facial features coordinates, such as the pupil centers, the inside and outside corners of eyes, and the widows peak in the hairline. The coordinates were used to calculate 20 individual distances, including the width of the mouth and of the eyes. A human could process about 40 pictures an hour, building a database of these computed distances. A computer would then automatically compare the distances for each photograph, calculate the difference between the distances, and return the closed records as a possible match.
In 1970, Takeo Kanade publicly demonstrated a face-matching system that located anatomical features such as the chin and calculated the distance 



\subsection{Scope Definition}
The project encompasses a dual-portal Web Application serving Students and Teachers. It includes self-service biometric facial registration, real-time facial recognition scanning, automated generation of CSV and PDF analytics reports, and an immutable audit logging system for manual interventions.

\subsection{Stakeholder Identification}
\begin{itemize}
    \item \textbf{Students}: Register facial biometrics and track personal attendance metrics.
    \item \textbf{Teachers}: Execute real-time facial attendance scanning, edit database records, and generate analytical reports.
    \item \textbf{Administrators}: Audit the immutable logs to track manual interventions in the database.
\end{itemize}

\subsection{Feasibility Study}
\subsubsection{Technical Feasibility}
Highly feasible. Web browsers natively support the MediaDevices API, and lightweight neural networks like \texttt{face-api.js} run entirely client-side via WebGL. FastAPI provides rapid, asynchronous backend communication.

\subsubsection{Economic Feasibility}
Highly feasible. The system leverages open-source technologies (Python, FastAPI, SQLite) and can be deployed on free-tier cloud environments like Render.

\subsubsection{Operational Feasibility}
Highly feasible. The UI is designed with a premium, accessible dark-mode aesthetic requiring minimal technical training.

\section{Requirement Engineering}
\subsection{Functional Requirements Specification}
\begin{enumerate}
    \item Authentication separating Student and Teacher workflows.
    \item Biometric Enrollment for students securely registering their facial vectors.
    \item Live video stream processing for Teachers to detect faces against the database.
    \item Dynamic generation of PDF and CSV Blacklist reports.
\end{enumerate}

\subsection{Non-Functional Requirements}
\begin{itemize}
    \item \textbf{Performance}: WebRTC must process frames at $\ge 15$ FPS.
    \item \textbf{Security}: Biometric data must be serialized 128D vectors.
    \item \textbf{Availability}: The FastAPI backend must handle 500+ concurrent connections.
\end{itemize}

\section{Software Development Life Cycle (SDLC) Planning}
\subsection{Selection of SDLC Model}
The Agile Methodology was selected to adapt to shifting stakeholder requirements, particularly during the iterative refinement of the reporting modules.
Agile software development is an umbrella term for approaches to developing software that reflect the values and principles agreed upon by The Agile Alliance, a group of 17 software practitioners, in 2001. As documented in their Manifesto for Agile Software Development, the practitioners value: 

Individuals and interactions over processes and tools
Working software over comprehensive documentation
Customer collaboration over contract negotiation
Responding to change over following a plan
The practitioners cite inspiration from new practices at the time including extreme programming, scrum, dynamic systems development method, adaptive software development, and being sympathetic to the need for an alternative to documentation-driven, heavyweight software development processes.
Many software development practices emerged from the agile mindset. These agile-based practices, sometimes called Agile (with a capital A), include requirements, discovery, and solutions improvement through the collaborative effort of self-organizing and cross-functional teams with their customer(s)/end user(s).
While there is much anecdotal evidence that the agile mindset and agile-based practices improve the software development process, the empirical evidence is limited and less than conclusive.


== History ==
Iterative and incremental software development methods can be traced back as early as 1957, with evolutionary project management and adaptive software development emerging in the early 1970s.
During the 1990s, a number of lightweight software development methods evolved in reaction to the prevailing heavyweight methods (often referred to collectively as waterfall) that critics described as overly regulated, planned, and micromanaged. These lightweight methods included rapid application development (RAD), from 1991; the unified process (UP) and dynamic systems development method (DSDM), both from 1994; Scrum, from 1995; Crystal Clear and extreme programming (XP), both from 1996; and feature-driven development (FDD), from 1997. Although these all originated before the publication of the Agile Manifesto, they are now collectively referred to as agile software development methods.
Already since 1991 similar changes had been underway in manufacturing and management thinking derived from lean management.
In 2001, seventeen software developers met at a resort in Snowbird, Utah to discuss lightweight development methods. They were Kent Beck (Extreme Programming), Ward Cunningham (Extreme Programming), Dave Thomas (Pragmatic Programming, Ruby), Jeff Sutherland (Scrum), Ken Schwaber (Scrum), Jim Highsmith (Adaptive Software Development), Alistair Cockburn (Crystal), Robert C. Martin (SOLID), Mike Beedle (Scrum), Arie van Bennekum, Martin Fowler (OOAD and UML), James Grenning, Andrew Hunt (Pragmatic Programming, Ruby), Ron Jeffries (Extreme Programming), Jon Kern, Brian Marick (Ruby, Test-driven development), and Steve Mellor (OOA). The group, The Agile Alliance, published the Manifesto for Agile Software Development.
In 2005, a group headed by Cockburn and Highsmith wrote an addendum of project management principles, the PM Declaration of Interdependence, to guide software project management according to agile software development methods.
In 2009, a group working with Martin wrote an extension of software development principles, the Software Craftsmanship Manifesto, to guide agile software development according to professional conduct and mastery.
In 2011, the Agi



\subsection{Gantt Chart Roadmap}
\begin{figure}[h]
    \centering
    \includegraphics[width=0.8\textwidth]{gantt_chart.png}
    \caption{Agile Development Roadmap}
\end{figure}

\section{System Modeling Using UML}
\subsection{Use Case Diagram}
\begin{figure}[h]
    \centering
    \includegraphics[width=0.6\textwidth]{use_case.png}
    \caption{Use Case Diagram mapping actor boundaries.}
\end{figure}

\subsection{Class Diagram}
\begin{figure}[h]
    \centering
    \includegraphics[width=0.6\textwidth]{class_diagram.png}
    \caption{Class Diagram illustrating SQLAlchemy ORM models.}
\end{figure}

\subsection{Sequence Diagram}
\begin{figure}[h]
    \centering
    \includegraphics[width=0.8\textwidth]{sequence_diagram.png}
    \caption{Sequence Diagram showing WebRTC video loop.}
\end{figure}

\subsection{Entity Relationship (ER) Diagram}
\begin{figure}[h]
    \centering
    \includegraphics[width=0.6\textwidth]{er_diagram.png}
    \caption{ER Diagram illustrating Relational keys.}
\end{figure}

\subsection{Activity Diagram}
\begin{figure}[h]
    \centering
    \includegraphics[width=0.4\textwidth]{activity_diagram.png}
    \caption{Activity Diagram tracking Audit Ledger modifications.}
\end{figure}

\chapter{Implementation, Testing, Deployment \& Evaluation Phase}

\section{Application Development}
\subsection{Frontend \& Backend Implementation}
The backend relies on the FastAPI asynchronous architecture, while the frontend utilizes Server-Side Rendered Jinja2 templates styled with TailwindCSS.

\subsection{Camera \& Face Detection Implementation (Code Snippet)}
Below is the critical implementation logic within the application development phase showing the WebRTC camera stream and face-api.js bounding box rendering.

\begin{lstlisting}[language=JavaScript, caption=WebRTC Face Detection Implementation]
async function detectLoop() {
    if(!running || video.paused || video.ended) return;

    const detections = await faceapi.detectAllFaces(video)
        .withFaceLandmarks()
        .withFaceDescriptors();
    
    const displaySize = { width: video.videoWidth, height: video.videoHeight };
    faceapi.matchDimensions(overlay, displaySize);
    const resizedDetections = faceapi.resizeResults(detections, displaySize);
    
    const ctx = overlay.getContext('2d');
    ctx.clearRect(0, 0, overlay.width, overlay.height);

    const results = resizedDetections.map(d => faceMatcher.findBestMatch(d.descriptor));

    results.forEach((result, i) => {
        const box = resizedDetections[i].detection.box;
        let text = result.label;
        let color = 'rgba(0, 240, 255, 1)'; 
        
        let drawBoxObject = box;
        // Invert X coordinate if the camera is mirrored via CSS
        if (currentFacingMode !== 'environment') {
            const newX = overlay.width - (box.x + box.width);
            drawBoxObject = new faceapi.Rect(newX, box.y, box.width, box.height);
        }

        const drawBox = new faceapi.draw.DrawBox(drawBoxObject, { 
            label: text,
            lineWidth: 2,
            boxColor: color
        });
        drawBox.draw(overlay);
    });
    requestAnimationFrame(detectLoop);
}
\end{lstlisting}

\subsection{Backend PDF Generation Implementation (Code Snippet)}
\begin{lstlisting}[language=Python, caption=Dynamic PDF Streaming Implementation]
from fpdf import FPDF
from fastapi import Response

def generate_pdf_table(title: str, headers: list, rows: list):
    pdf = FPDF(orientation='L' if len(headers) > 6 else 'P')
    pdf.add_page()
    pdf.set_font("helvetica", "B", 16)
    pdf.cell(0, 10, title, align="C")
    pdf.ln(15)
    
    col_width = pdf.epw / len(headers)
    
    pdf.set_font("helvetica", "B", 10)
    for header in headers:
        pdf.cell(col_width, 10, str(header), border=1, align="C")
    pdf.ln()
    
    pdf.set_font("helvetica", "", 10)
    for row in rows:
        for item in row:
            pdf.cell(col_width, 10, str(item), border=1, align="C")
        pdf.ln()
        
    return pdf.output()
\end{lstlisting}

\section{Integration \& System Testing}
\subsection{Testing Procedures}
Unit testing was utilized to handle edge cases, such as `ZeroDivisionError` during PDF calculations. Black-box testing evaluated the UI login constraints.
Software testing is the act of checking whether software meets its intended objectives and satisfies expectations.
Software testing can provide objective, independent information about the quality of software and the risk of its failure to a user or sponsor or any other stakeholder.
Software testing can determine the correctness of software for specific scenarios but cannot determine correctness for all scenarios. It cannot find all bugs.
Based on the criteria for measuring correctness from an oracle, software testing employs principles and mechanisms that might recognize a problem. Examples of oracles include specifications, contracts, comparable products, past versions of the same product, inferences about intended or expected purpose, user or customer expectations, relevant standards, and applicable laws.
Software testing can be functional or non-functional in nature.
Software testing is often dynamic in nature: running the software to verify actual output matches expected. It can also be static in nature: reviewing code and its associated documentation.
Software testing is often used to answer the question: Does the software do what it is supposed to do and what it needs to do?
Information learned from software testing may be used to improve the process by which software is developed.
A commonly suggested approach to automated testing is the "test pyramid," wherein most of the tests are unit tests, followed by a smaller set of integration tests and finally a few end-to-end (e2e) tests.


== Economics ==
A study conducted by NIST in 2002 reported that software bugs cost the U.S. economy $59.5 billion annually. More than a third of this cost could be avoided if better software testing was performed.
Outsourcing software testing because of costs is very common, with China, the Philippines, India and Pakistan being preferred destinations.


== History ==
Glenford J. Myers initially introduced the separation of debugging from testing in 1979. Although his attention was on breakage testing ("A successful test case is one that detects an as-yet undiscovered error."), it illustrated the desire of the software engineering community to separate fundamental development activities, such as debugging, from that of verification.
Software testing typically includes handling software bugs – a defect in the code that causes an undesirable result. Bugs generally slow testing progress and involve programmer assistance to debug and fix.
Not all defects cause a failure. For example, a defect in dead code will not be considered a failure.
A defect that does not cause failure at one point in time may lead to failure later due to environmental changes. Examples of environment change include running on new computer hardware, changes in data, and interacting with different software.


== Goals ==
Software testing is typically goal driven.


=== Finding bugs ===
Software testing typically includes handling software bugs – a defect in the code that causes an undesirable result. Bugs generally slow testing progress and involve programmer assistance to debug and fix.
Not all defects cause a failure. For example, a defect in dead code will not be considered a failure.
A defect that does not cause failure at one point in time may lead to failure later due to environmental changes. Examples of environment change include running on new computer hardware, changes in data, and interacting with different software.
A single defect may result in multiple failure symptoms



\section{Application Deployment}
\subsection{Cloud Hosting}
The application was deployed to the Render PaaS ecosystem, leveraging \texttt{requirements.txt} for continuous delivery pipelines.
Cloud computing is defined by the International Organization for Standardization (ISO) as "a paradigm for enabling network access to a scalable and elastic pool of shareable physical or virtual resources with self-service provisioning and administration on demand". It is commonly referred to as "the cloud".


== Characteristics ==
In 2011, the National Institute of Standards and Technology (NIST) identified five "essential characteristics" for cloud systems. Below are the exact definitions according to NIST:

On-demand self-service: "A consumer can unilaterally provision computing capabilities, such as server time and network storage, as needed automatically without requiring human interaction with each service provider."
Broad network access: "Capabilities are available over the network and accessed through standard mechanisms that promote use by heterogeneous thin or thick client platforms (e.g., mobile phones, tablets, laptops, and workstations)."
Resource pooling: " The provider's computing resources are pooled to serve multiple consumers using a multi-tenant model, with different physical and virtual resources dynamically assigned and reassigned according to consumer demand."
Rapid elasticity: "Capabilities can be elastically provisioned and released, in some cases automatically, to scale rapidly outward and inward commensurate with demand. To the consumer, the capabilities available for provisioning often appear unlimited and can be appropriated in any quantity at any time."
Measured service: "Cloud systems automatically control and optimize resource use by leveraging a metering capability at some level of abstraction appropriate to the type of service (e.g., storage, processing, bandwidth, and active user accounts). Resource usage can be monitored, controlled, and reported, providing transparency for both the provider and consumer of the utilized service.
By 2023, the International Organization for Standardization (ISO) had expanded and refined the list.


== History ==

The history of cloud computing extends to the 1960s, with the initial concepts of time-sharing becoming popularized via remote job entry (RJE). The "data center" model, where users submitted jobs to operators to run on mainframes, was predominantly used during this era. This period saw broad experimentation with making large-scale computing power more accessible through time-sharing, while optimizing infrastructure, platforms, and applications to improve efficiency for end users.
The "cloud" metaphor for virtualized services dates to 1994, when it was used by General Magic for the universe of "places" that mobile agents in the Telescript environment could "go". The metaphor is credited to David Hoffman, a General Magic communications specialist, based on its long-standing use in networking and telecom. The expression cloud computing became more widely known in 1996 when Compaq Computer Corporation drew up a business plan for future computing and the Internet. The company's ambition was to supercharge sales with "cloud computing-enabled applications". The business plan foresaw that online consumer file storage would likely be commercially successful. As a result, Compaq decided to sell server hardware to internet service providers.
In the 2000s, the application of cloud computing began to take shape with the establishment of Amazon Web Services (AWS) in 2002, which allowed developers to build applications independently. In 2006 Amazon Simple Storage Service, known



\section{Final Documentation \& Extensions}
\section{Theoretical Background: Machine learning}
Machine learning (ML) is a field of study in artificial intelligence concerned with the development and study of statistical algorithms that can learn from  data and generalize to unseen data, and thus perform tasks without being explicitly programmed.
Statistics and mathematical optimisation methods compose the foundations of machine learning. Data mining is a related field of study, focusing on exploratory data analysis (EDA) through unsupervised learning.
From a theoretical viewpoint, probably approximately correct learning provides a mathematical and statistical framework for describing machine learning. Most traditional machine learning and deep learning algorithms can be described as empirical risk minimisation under this framework.
Advances in the field of deep learning have allowed neural networks, a class of statistical algorithms, to surpass many previous machine learning approaches in performance. The resulting dominance of machine learning within the field of AI from the late 2000s and early 2010s onwards was such that the terms "machine learning" and "artificial intelligence" are now often used interchangeably, even though they are not synonymous.


== History ==

The term machine learning was coined in 1959 by Arthur Samuel, an IBM employee and pioneer in the field of computer gaming and artificial intelligence. The synonym self-teaching computers was also used during this time period.
The earliest machine learning program was introduced in the 1950s, when Samuel invented a computer program that calculated the chance of winning in checkers for each side, but the history of machine learning is rooted in decades of efforts to study human cognitive processes. In 1949, Canadian psychologist Donald Hebb published the book The Organization of Behavior, in which he introduced a theoretical neural structure formed by certain interactions among nerve cells. The Hebbian theory of neuron interaction set the groundwork for how many machine learning algorithms work, with connected artificial neurons changing the strength of their connections based on data. Other researchers who have studied human cognitive systems contributed to the modern machine learning technologies as well, including  Walter Pitts and Warren McCulloch, who proposed the first mathematical model of neural networks including algorithms that mirror human thought processes.
By the early 1960s, an experimental "learning machine" with punched tape memory, called Cybertron, had been developed by Raytheon Company to analyse sonar signals, electrocardiograms, and speech patterns using rudimentary reinforcement learning. It was repetitively "trained" by a human operator/teacher to recognise patterns and equipped with a "goof" button to cause it to reevaluate incorrect decisions. A representative book on research into machine learning during the 1960s was Nils Nilsson's book "Learning Machines", dealing mostly with machine learning for pattern classification. Interest related to pattern recognition continued into the 1970s, as described by Duda and Hart in 1973. In 1981, a report was given on using teaching strategies so that an artificial neural network learns to recognise 40 characters (26 letters, 10 digits, and 4 special symbols) from a computer terminal.
Tom M. Mitchell provided a widely quoted, more formal definition of the algorithms studied in the machine learning field: "A computer program is said to learn from experience E with respect to some class of tasks T and pe

\section{Theoretical Background: Artificial neural network}
In machine learning, a neural network (NN) or artificial neural network (ANN) is a computational model inspired by the structure and functions of biological neural networks.
A neural network consists of connected units or nodes called artificial neurons, which loosely model the neurons in the brain. These are connected by edges, which model the synapses in the brain. Each artificial neuron receives signals from connected neurons, then processes them and sends a signal to other connected neurons. The "signal" is a real number, and the output of each neuron is computed by some non-linear function of the totality of its inputs, called the activation function. The strength of the signal at each connection is determined by a weight, which adjusts as part of the training process.
Groups of neurons are aggregated into layers. Each layer performs a transformation on its inputs. Signals travel from the first layer (the input layer) to the last layer (the output layer), typically passing through multiple intermediate layers (hidden layers). A network is typically called a deep neural network if it has at least two hidden layers. Deep neural networks are capable of learning sophisticated hierarchical representations.
Training neural networks is a compute-intensive process, accelerated by  the use of graphics processing units (GPUs), and large datasets.

In reality, such textures and outlines would not be represented by single nodes, but rather by associated weight patterns of multiple nodes.
Architectural innovations such as convolutional neural networks (CNNs) significantly improved performance in computer vision tasks, while recurrent neural networks (RNNs) enabled modeling of sequential data such as speech and time-series information. Transformer architectures introduced attention mechanisms that allow neural networks to model long-range dependencies in data and have been the basis of large language models.
Artificial neural networks are used for myriad tasks including chatbots, large-scale text, image, and video generation, and robotics.


== History ==


=== Mathematical foundations ===
Deep neural networks are based on statistics developed over 200 years ago. The simplest kind of feedforward neural network (FNN) is a linear network, which consists of a single layer of output nodes with linear activation functions; the inputs are fed directly to the outputs via weights. The sum of the products of the weights and the inputs is calculated at each node. The mean squared errors between these calculated outputs and the given target values are minimized by adjusting to the weights. This technique is the method of least squares or linear regression. It was used to find a rough linear fit to a set of points by Legendre (1805) and Gauss (1795) for the prediction of planetary movement.


=== Perceptrons ===
Computers are based on John von Neumann's model. They execute explicit lists of instructions with access to memory to record their changing state. Neural networks instead originated from efforts to model information processing in biological systems via connectionism. Unlike the von Neumann model, connectionist computing does not separate memory and processing.
Warren McCulloch and Walter Pitts (1943) considered a non-learning computational model for neural networks. This model paved the way for research to split into one branch focused on biological processes and another focused on artificial intelligence. McCulloch and Pitts also developed mathemat

\section{Theoretical Background: Convolutional neural network}
A convolutional neural network (CNN) is a type of feedforward neural network that learns features via filter (or kernel) optimization. This type of deep learning network has been applied to process and make predictions from many different types of data including text, images and audio. CNNs are the de-facto standard in deep learning-based approaches to computer vision and image processing, and have only recently been replaced—in some cases—by newer architectures such as the transformer.
Vanishing gradients and exploding gradients, seen during backpropagation in earlier neural networks, are prevented by the regularization that comes from using shared weights over fewer connections. For example, for each neuron in the fully-connected layer, 10,000 weights would be required for processing an image sized 100 × 100 pixels. However, applying cascaded convolution (or cross-correlation) kernels, only 25 weights for each convolutional layer are required to process 5x5-sized tiles. Higher-layer features are extracted from wider context windows, compared to lower-layer features.
Some applications of CNNs include: 

image and video recognition,
recommender systems,
image classification,
image segmentation,
medical image analysis,
natural language processing,
brain–computer interfaces, and
financial time series.
CNNs are also known as shift invariant or space invariant artificial neural networks, based on the shared-weight architecture of the convolution kernels or filters that slide along input features and provide translation-equivariant responses known as feature maps. Counter-intuitively, most convolutional neural networks are not invariant to translation, due to the downsampling operation they apply to the input.
Feedforward neural networks are usually fully connected networks, that is, each neuron in one layer is connected to all neurons in the next layer. The "full connectivity" of these networks makes them prone to overfitting data. Typical ways of regularization, or preventing overfitting, include: penalizing parameters during training (such as weight decay) or trimming connectivity (skipped connections, dropout, etc.) Robust datasets also increase the probability that CNNs will learn the generalized principles that characterize a given dataset rather than the biases of a poorly-populated set.
Convolutional networks were inspired by biological processes in that the connectivity pattern between neurons resembles the organization of the animal visual cortex. Individual cortical neurons respond to stimuli only in a restricted region of the visual field known as the receptive field. The receptive fields of different neurons partially overlap such that they cover the entire visual field.
CNNs use relatively little pre-processing compared to other image classification algorithms. This means that the network learns to optimize the filters (or kernels) through automated learning, whereas in traditional algorithms these filters are hand-engineered. This simplifies and automates the process, enhancing efficiency and scalability overcoming human-intervention bottlenecks.


== Architecture ==

A convolutional neural network consists of an input layer, hidden layers and an output layer. In a convolutional neural network, the hidden layers include one or more layers that perform convolutions. Typically this includes a layer that performs a dot product of the convolution kernel with the layer's input matrix. This product is usually the Frobenius inner pro

\section{Theoretical Background: Data privacy}
Information privacy, also known as data privacy or data protection, is the relationship between the collection and dissemination of data, technology, the public expectation of privacy, contextual information norms, and the legal and political issues surrounding them. 
Various types of personal information often come under privacy concerns. Privacy concerns exist wherever personally identifiable information or other sensitive information is collected, stored, used, and finally destroyed or deleted – in digital form or otherwise. Improper or non-existent disclosure control can be the root cause for privacy issues. Informed consent mechanisms including privacy policies and dynamic consent are important in communicating to data subjects the different uses of their personally identifiable information.


== History ==
In 1890, legal scholars Samuel D. Warren II and Louis Brandeis published "The Right to Privacy", an article explaining "the right to be let alone," which helped establish a right to privacy in the United States. Their guidelines set the foundation for how privacy can be understood in terms of personal data, and influenced the development of the privacy laws of the United States. 
In the United States, the Privacy Act of 1974 represented an achievement in information privacy law. The law set a Code of Fair Information Practice to govern how federal agencies collect personal information and data. After the Watergate and COINTELPRO scandals involving illegal government use of surveillance, this privacy law was intended to help restore public trust in the government. The Act allowed a person to have the right to access information about themselves and know how their information would be used.
The growth of the internet's popularity in the early 2000s presented different privacy concerns as many companies and organizations began to collect user's personal data. Scholars such as Daniel Solove explained the importance of informed consent.
In 2018, the Facebook–Cambridge Analytica data scandal introduced new privacy risks to the public. A consulting firm called Cambridge Analytica collected personal user data from over 85 million Facebook users without informed consent. The data collected was used to assist in politically targeted ads during the 2016 U.S. presidential election. This resulted in the United States Federal Trade Commission fining Facebook $5 billion. The Cambridge Analytica scandal influenced new legislation such as the European Union's General Data Protection Regulation.


== Types of information privacy ==
Data privacy issues may arise in response to information from a wide range of sources, such as:

Healthcare records
Criminal justice investigations and proceedings
Financial institutions and transactions
Biological traits, such as genetic material
Residence and geographic records
Privacy breach
Location-based service and geolocation
Web surfing behavior or user preferences using persistent cookies
Academic research


=== Financial ===

Information about a person's financial transactions, including the amount of assets, positions held in stocks or funds, outstanding debts, and purchases can be sensitive. If criminals gain access to information such as a person's accounts or credit card numbers, that person could become the victim of fraud or identity theft. Information about a person's purchases can reveal a great deal about that person's history, such as places they have visited, whom they have contact with, products t

\section{Theoretical Background: Cryptography}
Cryptography, or cryptology, is the practice and study of techniques for secure communication in the presence of adversarial behavior. More generally, cryptography is about constructing and analyzing protocols that prevent third parties or the public from reading private messages. Modern cryptography exists at the intersection of the disciplines of mathematics, computer science, information security, electrical engineering, digital signal processing, physics, and others. Core concepts related to information security (data confidentiality, data integrity, authentication and non-repudiation) are also central to cryptography. Practical applications of cryptography include electronic commerce, chip-based payment cards, digital currencies, computer passwords and military communications.
Cryptography prior to the modern age was effectively synonymous with encryption, converting readable information (plaintext) to unintelligible nonsense text (ciphertext), which can only be read by reversing the process (decryption). The sender of an encrypted (coded) message shares the decryption (decoding) technique only with the intended recipients to preclude access from adversaries. The cryptography literature often uses the names "Alice" (or "A") for the sender, "Bob" (or "B") for the intended recipient, and "Eve" (or "E") for the eavesdropping adversary. Since the development of rotor cipher machines in World War I and the advent of computers in World War II, cryptography methods have become increasingly complex and their applications more varied.
Modern cryptography is heavily based on mathematical theory and computer science practice; cryptographic algorithms are designed around computational hardness assumptions, making such algorithms hard to break in actual practice by any adversary. While it is theoretically possible to break into a well-designed system, it is infeasible in actual practice to do so. Such schemes, if well designed, are therefore termed "computationally secure". Theoretical advances (e.g., improvements in integer factorization algorithms) and faster computing technology require these designs to be continually reevaluated and, if necessary, adapted. Information-theoretically secure schemes that provably cannot be broken even with unlimited computing power, such as the one-time pad, are much more difficult to use in practice than the best theoretically breakable but computationally secure schemes.
The growth of cryptographic technology has raised a number of legal issues in the Information Age. Cryptography's potential for use as a tool for espionage and sedition has led many governments to classify it as a weapon and to limit or even prohibit its use and export. In some jurisdictions where the use of cryptography is legal, laws permit investigators to compel the disclosure of encryption keys for documents relevant to an investigation. Cryptography also plays a major role in digital rights management and copyright infringement disputes with regard to digital media.


== Terminology ==

The first use of the term "cryptograph" (as opposed to "cryptogram") dates back to the 19th century – originating from "The Gold-Bug", a story by Edgar Allan Poe. The etymology of the term goes back to the Greek language; it is formed by two constituents: "crypton" (hidden), and "grapho" (to write).
Until modern times, cryptography referred almost exclusively to "encryption", which is the process of converting ordinary information (called plaintext) int

\section{Theoretical Background: Object-relational mapping}
Object–relational mapping (ORM, O/RM, and O/R mapping tool) in computer science is a programming technique for converting data between a relational database and the memory (usually the heap) of an object-oriented programming language. This creates, in effect, a virtual object database that can be used from within the program.
In object-oriented programming, data-management tasks act on objects that combine scalar values into objects. For example, consider an address book entry that represents a single person along with zero or more phone numbers and zero or more addresses. This could be modeled in an object-oriented implementation by a "Person object" with an attribute/field to hold each data item that the entry comprises: the person's name, a list of phone numbers, and a list of addresses. The list of phone numbers would itself contain "PhoneNumber objects" and so on. Each such address-book entry is treated as a single object by the programming language (it can be referenced by a single variable containing a pointer to the object, for instance). Various methods can be associated with the object, such as methods to return the preferred phone number, the home address, and so on.
By contrast, relational databases, such as SQL, group scalars into tuples, which are then enumerated in tables. Tuples and objects have some general similarity, in that they are both ways to collect values into named fields such that the whole collection can be manipulated as a single compound entity. They have many differences, though, in particular: lifecycle management (row insertion and deletion, versus garbage collection or reference counting), references to other entities (object references, versus foreign key references), and inheritance (non-existent in relational databases). As well, objects are managed on-heap and are under full control of a single process, while database tuples are shared and must incorporate locking, merging, and retry. Object–relational mapping provides automated support for mapping tuples to objects and back, while accounting for all of these differences.
The heart of the problem involves translating the logical representation of the objects into an atomized form that is capable of being stored in the database while preserving the properties of the objects and their relationships so that they can be reloaded as objects when needed. If this storage and retrieval functionality is implemented, the objects are said to be persistent.


== Overview ==
Implementation-specific details of storage drivers are generally wrapped in an API in the programming language in use, exposing methods to interact with the storage medium in a way which is simpler and more in line with the paradigms of surrounding code.
The following is a simple example, written in C\# code, to execute a query written in SQL using a database engine.

In contrast, the following makes use of an ORM-job API which makes it possible to write code that naturally makes use of the features of the language.

The case above makes use of an object representing the storage repository and methods of that object. Other frameworks might provide code as static methods, as in the example below, and yet other methods may not implement an object-oriented system at all. Often the choice of paradigm is made for the best fit of the ORM into the surrounding language's design principles.


== Comparison with traditional data access techniques ==
Compared to traditional techniques of exchange betwee

\section{Theoretical Background: Single-page application}
A single-page application (SPA) is a web application or website that interacts with the user by dynamically rewriting the current web page with new data from the web server, instead of the default method of loading entire new pages. The goal is faster transitions that make the website feel more like a native app.
In a SPA, a page refresh never occurs; instead, all necessary HTML, JavaScript, and CSS code is either retrieved by the browser with a single page load, or the appropriate resources are dynamically loaded and added to the page as necessary, usually in response to user actions.


== History ==
The origins of the term single-page application are unclear, though the concept was discussed at least as early as 2003 by technology evangelists from Netscape. Stuart Morris, a programming student at Cardiff University, Wales, wrote the self-contained website at slashdotslash.com with the same goals and functions in April 2002, and later the same year Lucas Birdeau, Kevin Hakman, Michael Peachey and Clifford Yeh described a single-page application implementation in US patent 8,136,109. Earlier forms were called rich web applications.
JavaScript can be used in a web browser to display the user interface (UI), run application logic, and communicate with a web server. Mature free libraries are available that support the building of a SPA, reducing the amount of JavaScript code developers have to write.


== Technical approaches ==
There are various techniques available that enable the browser to retain a single page even when the application requires server communication.


=== Document hashes ===
HTML authors can leverage element IDs to show or hide different sections of the HTML document.  Then, using CSS, authors can use the :target pseudo-class selector to only show the section of the page which the browser navigated to.


=== JavaScript frameworks ===
Web browser JavaScript frameworks and libraries, such as Angular, Ember.js, ExtJS, Knockout.js, Meteor.js, React, Vue.js, and Svelte have adopted SPA principles. Aside from ExtJS, all of these are free. 

AngularJS is a discontinued fully client-side framework. AngularJS's templating is based on bidirectional UI data binding. Data-binding is an automatic way of updating the view whenever the model changes, as well as updating the model whenever the view changes. The HTML template is compiled in the browser. The compilation step creates pure HTML, which the browser re-renders into the live view. The step is repeated for subsequent page views. In traditional server-side HTML programming, concepts such as controller and model interact within a server process to produce new HTML views. In the AngularJS framework, the controller and model states are maintained within the client browser. Therefore, new pages are capable of being generated without any interaction with a server.
Angular 2+ is a SPA Framework developed by Google after AngularJS. There is a strong community of developers using this framework. The framework is updated twice every year. New features and fixes are frequently added in this framework.
Ember.js is a client-side JavaScript web application framework based on the model–view–controller (MVC) software architectural pattern. It allows developers to create scalable single-page applications by incorporating common idioms and best practices into a framework that provides a rich object model, declarative two-way data binding, computed properties, automatically updating templates po

\section{Theoretical Background: Web application}
A web application (or web app) is application software that is created with web technologies and runs via a web browser. Web applications emerged during the late 1990s and allowed the server to dynamically build a response to a request, in contrast to static web pages. 
Web applications are commonly distributed via a web server. There are several different tier systems that web applications use to communicate between the web browsers, the client interface, and server data. Each system has its own uses as they function in different ways. However, there are many security risks that developers must be aware of during development; proper measures to protect user data are vital.
Web applications are often constructed with the use of a web application framework. Single-page applications (SPAs) and progressive web apps (PWAs) are two architectural approaches to creating web applications that provide a user experience similar to native apps, including features such as smooth navigation, offline support, and faster interactions.
Web applications are often fully hosted on remote cloud services, can require a constant connection to them, and can replace conventional desktop applications for operating systems such as Microsoft Windows, thus facilitating the operation of software as a service as it grants the developer the power to tightly control billing based on use of the remote services as well as vendor lock-in by hosting data remotely. Modern browsers such as Chrome offer sandboxing for every browser tab which improves security and restricts access to local resources. No software installation is required as the app runs within the browser, which reduces the need for managing software installations. With the use of remote cloud services, customers do not need to manage servers, since that can be left to the developer and the cloud service, and can use the software with a relatively low power, low-resource PC such as a thin client. The source code of the application can stay the same across operating systems and devices of users with the use of responsive web design, since it only needs to be compatible with web browsers which adhere to web standards, making the code highly portable and saving on development time. Numerous JavaScript frameworks and CSS frameworks facilitate development.


== History ==
The concept of a "web application" was first introduced in the Java language in the Servlet Specification version 2.2, which was released in 1999. At that time, both JavaScript and XML had already been developed, but the XMLHttpRequest object had only been recently introduced on Internet Explorer 5 as an ActiveX object. Beginning around the early 2000s, applications such as "Myspace (2003), Gmail (2004), Digg (2004), [and] Google Maps (2005)," started to make their client sides more and more interactive. A web page script is able to contact the server for storing and retrieving data without downloading an entire web page. The practice became known as Ajax in 2005. Eventually this was replaced by web APIs using JSON, accessed via JavaScript asynchronously on the client side.
In earlier computing models like client-server, the processing load for the application was shared between code on the server and code installed on each client locally. In other words, an application had its own pre-compiled client program which served as its user interface and had to be separately installed on each user's personal computer. An upgrade to the server-side code o

\section{Theoretical Background: Computer vision}
Computer vision tasks include methods for acquiring, processing, analyzing, and understanding digital images, and extraction of high-dimensional data from the real world in order to produce numerical or symbolic information, e.g. in the form of decisions. "Understanding" in this context signifies the transformation of visual images into descriptions of the world that make sense to thought processes and can elicit appropriate action. This image understanding can be seen as the disentangling of symbolic information from image data using models constructed with the aid of geometry, physics, statistics, and learning theory.
The scientific discipline of computer vision is concerned with the theory behind artificial systems that extract information from images. Image data can take many forms, such as video sequences, views from multiple cameras, multi-dimensional data from a 3D scanner, 3D point clouds from LiDaR sensors, or medical scanning devices. The technological discipline of computer vision seeks to apply its theories and models to the construction of computer vision systems.
Subdisciplines of computer vision include scene reconstruction, object detection, event detection, activity recognition, video tracking, object recognition, 3D pose estimation, learning, indexing, motion estimation, visual servoing, 3D scene modeling, and image restoration.


== Definition ==
Computer vision is an interdisciplinary field that deals with how computers can be made to gain high-level understanding from digital images or videos. From the perspective of engineering, it seeks to automate tasks that the human visual system can do. "Computer vision is concerned with the automatic extraction, analysis, and understanding of useful information from a single image or a sequence of images. It involves the development of a theoretical and algorithmic basis to achieve automatic visual understanding." As a scientific discipline, computer vision is concerned with the theory behind artificial systems that extract information from images. The image data can take many forms, such as video sequences, views from multiple cameras, or multi-dimensional data from a medical scanner. As a technological discipline, computer vision seeks to apply its theories and models for the construction of computer vision systems. Machine vision refers to a systems engineering discipline, especially in the context of factory automation. In more recent times, the terms computer vision and machine vision have converged to a greater degree.


== History ==
In the late 1960s, computer vision began at universities that were pioneering artificial intelligence. It was meant to mimic the human visual system as a stepping stone to endowing robots with intelligent behavior. In 1966, it was believed that this could be achieved through an undergraduate summer project, by attaching a camera to a computer and having it "describe what it saw".
What distinguished computer vision from the prevalent field of digital image processing at that time was a desire to extract three-dimensional structure from images with the goal of achieving full scene understanding. Studies in the 1970s formed the early foundations for many of the computer vision algorithms that exist today, including extraction of edges from images, labeling of lines, non-polyhedral and polyhedral modeling, representation of objects as interconnections of smaller structures, optical flow, and motion estimation.
The next decade saw studies based

\section{Theoretical Background: Algorithm}
In mathematics and computer science, an algorithm ( ) is any well-defined set of instructions that when followed terminates after a finite number of steps that comprise a solution to a given computational problem. Advanced algorithms may utilize loops and involve many conditionals that decide the next step based on the inputs provided, resulting in long sequences of steps before halting, but all algorithms terminate by definition.
In contrast to algorithms, heuristics might be applied to problems for which it is not possible for an algorithmic solution to exist, e.g., because the problem is undecidable or there is otherwise no way to define the correct or optimal result. For example, although social media recommender systems are erroneously called "algorithms" in popular media, they are accurately classified as heuristics, because there is no objective criteria by which a recommendation produced by such a system can be judged as "correct".
As an effective method, an algorithm can be expressed within both a finite amount of space and time and in a well-defined formal language for calculating a function. Starting from an initial state and input, a computation occurs at each step, eventually producing output and terminating.  Randomized algorithms incorporate random inputs which can be used to simulate non-deterministic output even though the transitions between states are deterministic.


== Etymology ==
Around 825 AD, Persian scientist and polymath al-Khwarizmi wrote kitāb al-ḥisāb al-hindī ("Book of Indian computation") and kitab al-jam' wa'l-tafriq al-ḥisāb al-hindī ("Addition and subtraction in Indian arithmetic"). In the early 12th century, Latin translations of these texts involving the Hindu–Arabic numeral system and arithmetic appeared, for example Liber Alghoarismi de practica arismetrice, attributed to John of Seville, and Liber Algoritmi de numero Indorum, attributed to Adelard of Bath. Here, alghoarismi or algoritmi is the Latinization of Al-Khwarizmi's name; the text starts with the phrase Dixit Algoritmi, or "Thus spoke Al-Khwarizmi".

 
The word algorism in English came to mean the use of place-value notation in calculations; it occurs in the Ancrene Wisse from circa 1225. By the time Geoffrey Chaucer wrote The Canterbury Tales in the late 14th century, he used a variant of the same word in describing augrym stones, stones used for place-value calculation. In the 15th century, under the influence of the Greek word ἀριθμός (arithmos, 'number'; cf. "arithmetic"), the Latin word was altered to algorithmus. By 1596, this form of the word was used in English, as algorithm, by Thomas Hood.


== Definition ==

One informal definition is "a set of rules that precisely defines a sequence of operations", which would include all computer programs, and any bureaucratic procedure
or cook-book recipe. In general, a program is an algorithm only if it stops eventually. Formally, algorithm is an explicit set of instructions to produce an output, that can be followed by a computer or a human performing specific operations on symbols.


== History ==


=== Ancient algorithms ===
Step-by-step procedures for solving mathematical problems have been recorded since antiquity. This includes in Babylonian mathematics (around 2500 BC), Egyptian mathematics (around 1550 BC), Indian mathematics (around 800 BC and later), the Ifa Oracle (around 500 BC), Greek mathematics (around 240 BC), Chinese mathematics (around 200 BC and later), and Arabic mathemat

\section{Theoretical Background: Database normalization}
Database normalization is the process of structuring a relational database in accordance with a series of normal forms to reduce data redundancy and improve data integrity. It was first proposed by British computer scientist Edgar F. Codd as part of his relational model.
Normalization entails organizing the columns (attributes) and tables (relations) of a database to ensure that their dependencies are properly enforced by database integrity constraints. It is accomplished by applying some formal rules either by a process of synthesis (creating a new database design) or decomposition (improving an existing database design).


== Objectives ==
A basic objective of the first normal form defined by Codd in 1970 was to permit data to be queried and manipulated using a "universal data sub-language" grounded in first-order logic. An example of such a language is SQL, though it is one that Codd regarded as seriously flawed.
The objectives of normalization beyond 1NF (first normal form) were stated by Codd as:

To free the collection of relations from undesirable insertion, update and deletion dependencies.
To reduce the need for restructuring the collection of relations, as new types of data are introduced, and thus increase the life span of application programs.
To make the relational model more informative to users.
To make the collection of relations neutral to the query statistics, where these statistics are liable to change as time goes by.

When an attempt is made to modify (update, insert into, or delete from) a relation, the following undesirable side effects may arise in relations that have not been sufficiently normalized:

Insertion anomaly
There are circumstances in which certain facts cannot be recorded at all. For example, each record in a "Faculty and Their Courses" relation might contain a Faculty ID, Faculty Name, Faculty Hire Date, and Course Code. Therefore, the details of any faculty member who teaches at least one course can be recorded, but a newly hired faculty member who has not yet been assigned to teach any courses cannot be recorded, except by setting the Course Code to null.
Update anomaly
The same information can be expressed on multiple rows; therefore updates to the relation may result in logical inconsistencies. For example, each record in an "Employees' Skills" relation might contain an Employee ID, Employee Address, and Skill; thus a change of address for a particular employee may need to be applied to multiple records (one for each skill). If the update is only partially successful – the employee's address is updated on some records but not others – then the relation is left in an inconsistent state. Specifically, the relation provides conflicting answers to the question of what this particular employee's address is.
Deletion anomaly
Under certain circumstances, the deletion of data representing certain facts necessitates the deletion of data representing completely different facts. The "Faculty and Their Courses" relation described in the previous example suffers from this type of anomaly, for if a faculty member temporarily ceases to be assigned to any courses, the last of the records on which that faculty member appears must be deleted, effectively also deleting the faculty member, unless the Course Code field is set to null.


=== Minimize redesign when extending the database structure ===
A fully normalized database can be extended to accommodate new types of data with minimal changes to its existing stru


\end{document}
